\documentclass[12pt]{article}
\usepackage[utf8]{inputenc}
\usepackage{amsmath,amssymb,amsfonts}
\usepackage{tikz}
\usetikzlibrary{arrows.meta, angles, quotes}
\usepackage[a4paper, margin=0.8in]{geometry}
\usepackage{enumitem}
\usepackage{mathptmx}
\usepackage{needspace}

\setlength{\parindent}{0pt}
\setlist[enumerate]{leftmargin=*, align=left}


\newcommand{\examanswerslot}[3]{%
  \par\nopagebreak\vspace*{#1}%
  \par\nopagebreak\vspace*{0.2cm}%
  \phantom{.} \hfill \makebox[#2\linewidth]{\dotfill}\hspace{4pt}\makebox[15pt][r]{{[#3]}}\par
}


\pagestyle{empty}

\begin{document}
\setlength{\parindent}{0pt}

% --- START OF QUESTION 1: E1.1_types_of_number ---
\needspace{7cm}
\noindent\textbf{Question 1} \par\vspace{0.2cm}
\begin{enumerate}[label=\textbf{(\alph*)}]
\item Write 540 as a product of its prime factors.
\examanswerslot{2.5cm}{0.2}{3}
\item Find the lowest common multiple (LCM) of 540 and 210.
\examanswerslot{2.5cm}{0.2}{3}
\par\vspace{1cm}

% --- START OF QUESTION 2: E1.2_sets_notation_venn ---
\needspace{7cm}
\noindent\textbf{Question 2} \par\vspace{0.2cm}
\begin{enumerate}[label=\textbf{(\alph*)}]
\item Shade the region

\par\vspace{0.15cm}\hspace*{0.6cm}$$\par\vspace{0.15cm}
A \cap B' \\
[MATH_END] \\
on the Venn diagram. \\
\examanswerslot{2.5cm}{0.2}{3}
\item The total number of elements in the universal set is

\par\vspace{0.15cm}\hspace*{0.6cm}$$\par\vspace{0.15cm}
\text{n}(\mathscr{E}) = 40
[MATH_END] \\
Given that \\

\par\vspace{0.15cm}\hspace*{0.6cm}$$\par\vspace{0.15cm}
\text{n}(A) = 18
[MATH_END] \\
, \\

\par\vspace{0.15cm}\hspace*{0.6cm}$$\par\vspace{0.15cm}
\text{n}(B) = 22
[MATH_END] \\
and \\

\par\vspace{0.15cm}\hspace*{0.6cm}$$\par\vspace{0.15cm}
\text{n}(A \cup B)' = 5
[MATH_END] \\
, find \\

\par\vspace{0.15cm}\hspace*{0.6cm}$$\par\vspace{0.15cm}
\text{n}(A \cap B)
[MATH_END] \\
. \\
\examanswerslot{2.5cm}{0.2}{3}
\par\vspace{1cm}

% --- START OF QUESTION 3: E1.3_powers_and_roots ---
\needspace{7cm}
\noindent\textbf{Question 3} \par\vspace{0.2cm}
\begin{enumerate}[label=\textbf{(\alph*)}]
\item Calculate

\par\vspace{0.15cm}\hspace*{0.6cm}$$\par\vspace{0.15cm}
\frac{\sqrt[3]{343} \times 2^{-3}}{5^0}
[MATH_END] \\
and write your answer as a fraction. \\
\examanswerslot{2.5cm}{0.2}{3}
\item Simplify

\par\vspace{0.15cm}\hspace*{0.6cm}$$\par\vspace{0.15cm}
\left( \frac{64x^6}{y^{-3}} \right)^{-\frac{2}{3}}
[MATH_END] \\
. \\
\examanswerslot{2.5cm}{0.2}{3}
\par\vspace{1cm}

% --- START OF QUESTION 4: E1.4_fractions_decimals_percentages ---
\needspace{7cm}
\noindent\textbf{Question 4} \par\vspace{0.2cm}
\begin{enumerate}[label=\textbf{(\alph*)}]
\item Write the recurring decimal

\par\vspace{0.15cm}\hspace*{0.6cm}$$\par\vspace{0.15cm}
0.2\dot{3}\dot{6} \\
[MATH_END] \\
as a fraction in its lowest terms. You must show all your working. \\
\examanswerslot{2.5cm}{0.2}{3}
\item Nitin scores 64 marks out of 80 in a chemistry test. In a physics test, he scores 78\% . Find the test in which Nitin achieved the higher percentage and state the difference between his two percentage scores.
\examanswerslot{2.5cm}{0.2}{3}
\par\vspace{1cm}

% --- START OF QUESTION 5: E1.5_ordering_magnitude ---
\needspace{7cm}
\noindent\textbf{Question 5} \par\vspace{0.2cm}
\begin{enumerate}[label=\textbf{(\alph*)}]
\item Write these quantities in order of size, starting with the smallest.

\par\vspace{0.15cm}\hspace*{0.6cm}$$\par\vspace{0.15cm}
\frac{22}{7} \quad\quad \pi \quad\quad 3.141 \quad\quad \sqrt{9.86}
[MATH_END] \\
\examanswerslot{2.5cm}{0.2}{3}
\item Write these numbers in order of size, starting with the smallest.

\par\vspace{0.15cm}\hspace*{0.6cm}$$\par\vspace{0.15cm}
2^{60} \quad\quad 3^{48} \quad\quad 5^{36} \quad\quad 6^{24} \\
[MATH_END] \\
\examanswerslot{2.5cm}{0.2}{3}
\par\vspace{1cm}
\end{document}
